\documentclass{../assignment}

\عنوان{\متن‌لاتین{Git}}
\ترم{پاییز ۱۴۰۵}
\مهلت{۲۸ مهر ۱۴۰۵}

\begin{document}

\عنوان‌ساز

\فهرست‌مطالب

\قسمت{مقدمه}

هدف این تمرین آشنایی با مفاهیم و تصمیم‌گیری در موقعیت‌های رایج کار با
\متن‌لاتین{Git} است. برای انجام تمرین یک هفته فرصت دارید.
همهٔ سؤال‌ها به‌صورت نوشتاری پاسخ داده می‌شوند. در سؤال‌های سناریومحور،
مراحل پیشنهادی، دستورهای لازم و دلیل انتخاب روش خود را توضیح دهید؛
اجرای دستورها، ساخت مخزن یا ارائهٔ تصویر و خروجی واقعی لازم نیست.
مجموع نمرهٔ تمرین ۲۰ است.

\قسمت{سؤال‌های مفهومی}

\زیرقسمت{سؤال یک: \متن‌لاتین{merge} و \متن‌لاتین{rebase}}
\نمره{۳}

دستورهای \متن‌لاتین{git merge} و \متن‌لاتین{git rebase} چه کاری انجام می‌دهند؟
تفاوت آن‌ها را از نظر نحوهٔ ترکیب تغییرات و شکل تاریخچهٔ کامیت‌ها توضیح دهید.
برای هر کدام یک موقعیت مناسب استفاده بیان کنید و توضیح دهید چرا در آن موقعیت
آن روش را انتخاب می‌کنید.

\زیرقسمت{سؤال دو: انتخاب بخشی از تغییرات}
\نمره{۳}

منظور از بخش تغییرات (\متن‌لاتین{hunk}) در خروجی مقایسهٔ تغییرات
(\متن‌لاتین{diff}) در \متن‌لاتین{Git} چیست؟
توضیح دهید چگونه می‌توان با استفاده از خط فرمان، تغییرات یک فایل را به‌صورت
تعاملی بررسی کرد و فقط بعضی از بخش‌ها را برای کامیت بعدی در ناحیهٔ
\متن‌لاتین{staging} قرار داد. دستورها و انتخاب‌های لازم را ذکر کنید و بگویید
این کار چه فایده‌ای برای ساختن کامیت‌های منسجم دارد.

\زیرقسمت{سؤال سوم: یافتن منشأ خطا}
\نمره{۳}

دستور \متن‌لاتین{git bisect} چگونه کار می‌کند؟ مراحل پیدا کردن کامیتی را که
یک خطا را وارد کرده است، از مشخص کردن یک کامیت سالم و یک کامیت معیوب تا پایان
جست‌وجو، توضیح دهید. یک مثال از کاربرد آن بیاورید و توضیح دهید چرا این روش
برای یافتن منشأ خطا مفید است.

\زیرقسمت{سؤال چهارم: بررسی تغییرات پیش از کامیت}
\نمره{۳}

فرض کنید فایل \متن‌لاتین{README.md} در مخزن \متن‌لاتین{tracked} است و قبلاً
کامیت شده است. بخشی از تغییرات جدید
آن را وارد ناحیهٔ \متن‌لاتین{staging} کرده‌اید و سپس، بدون اجرای دوبارهٔ
دستور افزودن به این ناحیه، همان فایل را مجدداً ویرایش کرده‌اید.
همچنین پوشه‌ای به نام \متن‌لاتین{tmp/} دارید که فقط شامل فایل‌های موقت و
\متن‌لاتین{untracked} است.

\شروع{شمارش}
\فقره با چه دستورهایی وضعیت مخزن، تغییرات هنوز واردنشده به ناحیهٔ
\متن‌لاتین{staging} و تغییرات آمادهٔ کامیت را بررسی می‌کنید؟
تفاوت دو دستهٔ تغییرات اخیر را توضیح دهید.
\فقره اگر در همین وضعیت یک کامیت معمولی و بدون گزینهٔ
\متن‌لاتین{-a} ثبت کنید، کدام تغییرات فایل در آن قرار می‌گیرند و کدام
تغییرات باقی می‌مانند؟ دلیل پاسخ خود را بیان کنید.
\فقره برای جلوگیری از اضافه‌شدن فایل‌های پوشهٔ \متن‌لاتین{tmp/} به کامیت‌های
بعدی، چه تنظیمی و در کدام فایل انجام می‌دهید؟ آیا همین تنظیم برای فایلی که
قبلاً \متن‌لاتین{tracked} شده است نیز کافی است؟ توضیح دهید.
\پایان{شمارش}

\زیرقسمت{سؤال پنجم: رفع \متن‌لاتین{merge conflict} در کار گروهی}
\نمره{۴}

شما و هم‌تیمی‌تان از یک کامیت مشترک شروع کرده‌اید. شما در \متن‌لاتین{branch}
\متن‌لاتین{feature-description} و هم‌تیمی‌تان در \متن‌لاتین{branch main}،
یک خط از فایل \متن‌لاتین{README.md} را به دو متن متفاوت تغییر داده و
کامیت کرده‌اید. همهٔ تغییرات ثبت شده‌اند و \متن‌لاتین{working tree} در وضعیت
\متن‌لاتین{clean} است؛ یعنی تغییر ثبت‌نشده‌ای ندارید.
هنگام \متن‌لاتین{merge} کردن \متن‌لاتین{branch} شما در \متن‌لاتین{main}،
\متن‌لاتین{Git} برای این خط \متن‌لاتین{conflict} گزارش می‌کند.

\شروع{شمارش}
\فقره چرا این \متن‌لاتین{conflict} رخ داده است؟ آیا هر تغییری که دو نفر در یک فایل
ایجاد کنند الزاماً باعث \متن‌لاتین{conflict} می‌شود؟
\فقره چگونه فایل‌های دارای \متن‌لاتین{conflict} و دو نسخهٔ متن را شناسایی می‌کنید؟
نقش \متن‌لاتین{conflict markers} در فایل را توضیح دهید.
\فقره فرض کنید اطلاعات هر دو تغییر باید حفظ شود. از زمان مشاهدهٔ \متن‌لاتین{conflict}
تا تکمیل \متن‌لاتین{merge} چه مراحلی را طی می‌کنید و چه دستورهایی به کار می‌برید؟
چگونه مطمئن می‌شوید که \متن‌لاتین{conflict markers} در متن نهایی باقی نمانده‌اند؟
\فقره اگر هنگام بررسی متوجه شوید که فعلاً نباید این \متن‌لاتین{merge} انجام شود،
با چه دستوری عملیات \متن‌لاتین{merge} در حال انجام را لغو می‌کنید؟
\پایان{شمارش}

\زیرقسمت{سؤال ششم: انتخاب روش مناسب برای بازگرداندن تغییرات}
\نمره{۴}

در هر یک از موقعیت‌های مستقل زیر، توضیح دهید چه می‌کنید، چه دستورهایی
به کار می‌برید و چه اثری بر \متن‌لاتین{working tree}، ناحیهٔ \متن‌لاتین{staging} یا
تاریخچهٔ کامیت‌ها می‌گذارند. فرض کنید هنگام اجرای دستورهای بازگردانی
\متن‌لاتین{conflict} رخ نمی‌دهد.

\شروع{شمارش}
\فقره در فایل \متن‌لاتین{tracked} به نام \متن‌لاتین{notes.txt} تغییری ایجاد کرده‌اید که
دیگر آن را نمی‌خواهید. این تغییر وارد ناحیهٔ \متن‌لاتین{staging} نشده است
و نسخهٔ موجود در این ناحیه با آخرین کامیت یکسان است. می‌خواهید فقط تغییرات
این فایل کنار گذاشته شوند و فایل‌های دیگر دست‌نخورده بمانند.
\فقره تغییرات \متن‌لاتین{notes.txt} را اشتباهاً وارد ناحیهٔ
\متن‌لاتین{staging} کرده‌اید، اما هنوز کامیت نکرده‌اید. می‌خواهید فایل از
این ناحیه خارج شود، ولی ویرایش‌های شما در \متن‌لاتین{working tree} حفظ شوند.
\فقره یک کامیت اشتباه معمولی، نه \متن‌لاتین{merge commit}، را روی
\متن‌لاتین{branch} مشترک \متن‌لاتین{main}، \متن‌لاتین{push} کرده‌اید و
هم‌تیمی‌ها نیز آن را دریافت کرده‌اند.
\متن‌لاتین{working tree} در وضعیت \متن‌لاتین{clean} است. می‌خواهید اثر این
کامیت را بدون حذف یا بازنویسی تاریخچه بازگردانید. روش مناسب چیست و چرا حذف
کامیت از تاریخچه و \متن‌لاتین{force push} در این موقعیت انتخاب مناسبی نیست؟
\پایان{شمارش}

\قسمت{نحوهٔ تحویل}

پاسخ هر شش سؤال را در یک گزارش، با شمارهٔ سؤال مربوط، تحویل دهید.
در پاسخ‌ها صرفاً نام دستور را ننویسید؛ گزینه‌ها و آرگومان‌های لازم،
دلیل انتخاب و نتیجهٔ مورد انتظار را نیز توضیح دهید. برای نام شاخه، مسیر
فایل یا شناسهٔ کامیت می‌توانید از مقادیر فرضی مشخص استفاده کنید.

\قوانین
\پایان‌ساز

\end{document}
